\begin{theorem}[Classification of the allowed zero-cycles on an abelian variety]
\label{c:abelian-point-classification}
\status{Proved over every algebraically closed field}
Let $A/F$ be an abelian variety over an algebraically closed field of any characteristic, and let $o\in\CH_0(A)_\Q$ have degree one. The following are equivalent:
\begin{enumerate}[label=(\roman*)]
\item all the cohomological point-and-diagonal relations hold in $T_k^{\CH}(A,o)$ for every $k$;
\item $\Delta_*o=o\boxtimes o$ in $\CH_0(A\times A)_\Q$;
\item $o=[a]$ in $\CH_0(A)_\Q$ for a point $a\in A(F)$.
\end{enumerate}
The point is unique modulo torsion in $A(F)$. In particular, specifying the origin is sufficient but is not necessary for the requested lifting.
\end{theorem}

\begin{proof}
The equality in (ii) is the cohomological relation $\delta_{12}y_1=y_1y_2$, since both classes are degree-one zero-cycles on $A^2$. Thus (i) implies (ii). If (iii) holds, simultaneous translation by $a$ on every factor preserves all partial diagonals and carries the origin class to $[a]$. Theorem~\ref{c:chow-positive} proves (i).

We prove the remaining implication algebraically, so that no characteristic-zero restriction is hidden. Put $g=\dim A$ and write $\widehat A$ for the dual abelian variety. For the normalized Poincar\'e bundle $\mathcal P$ on $A\times\widehat A$, define the rational Chow Fourier transform by
\[
 \mathcal F_A(z)=(p_2)_*\bigl(p_1^*z\,\exp(c_1(\mathcal P))\bigr).
\]
It is an isomorphism and $\mathcal F_{\widehat A}\mathcal F_A=(-1)^g[-1]^*$. These identities hold over the arbitrary algebraically closed field of \citet[Theorem 2.2, p.~156, with the field convention on p.~154]{c_mukai1981}: apply Grothendieck--Riemann--Roch to his Fourier--Mukai identity, using the trivial tangent bundle of an abelian variety. The Chern character identifies rational $K$-theory with rational Chow groups. We will also use the directly checked Poincar\'e identities below; over $\C$ they are formulas F1--F3 and Proposition 1 of \citet[pp.~647--648]{c_beauville1986}.

Set $z=\mathcal F_A(o)$. Its codimension-zero part is $\deg(o)=1$. Fourier transformation of (ii), using the defining Poincar\'e bundle and the projection formula, gives
\[
 m_{\widehat A}^*z=p_1^*z\,p_2^*z.
\]
Indeed Fourier transformation changes $\Delta_*$ into pullback by the dual homomorphism $m_{\widehat A}$, and changes an exterior product into the exterior product of the transforms. The ideal of positive-codimension classes is nilpotent, so the finite series $D=\log z$ is defined and satisfies
\[
 m_{\widehat A}^*D=p_1^*D+p_2^*D,
 \qquad [n]^*D=nD\quad(n\geq1).
\]

Write $D=\sum_{r=1}^gD_r$ by codimension. We verify that $D_r=0$ for $r\geq2$ without assuming the conjectural vanishing of any negative Beauville piece. Choose an integer $n\geq2$ prime to the characteristic. If $w_q$ is the codimension-$q$ part of $\mathcal F_{\widehat A}(D_r)$, its defining formula uses the power $c_1(\mathcal P)^{g-r+q}$. The identity $(1\times[n])^*\mathcal P=\mathcal P^{\otimes n}$ and flat base change therefore give
\[
 [n]^*w_q=n^{g-r+q}w_q.
\]
Fourier functoriality for the isogeny $[n]$, applied to $[n]^*D_r=nD_r$, gives $[n]_*w_q=nw_q$. Since $[n]_*[n]^*=n^{2g}$, a nonzero $w_q$ requires
\[
 n^{2g}=n^{g-r+q+1},\qquad q=g+r-1.
\]
The inequality $q\leq g$ forces $r\leq1$. Fourier inversion now gives $D=D_1$. On the N\'eron--Severi group $[n]^*$ acts by $n^2$, whereas it acts on $D_1$ by $n$. Hence $D_1\in\Pic^0(\widehat A)\otimes\Q$.

Biduality identifies $\Pic^0(\widehat A)$ with $A(F)$. This group is divisible: multiplication by every positive integer is a surjective isogeny and is surjective on points over an algebraically closed field, also when the integer is divisible by the characteristic. Thus there is $a\in A(F)$ with $D_1=c_1(\mathcal P_a)$ in rational Chow groups. Directly from its definition, $\mathcal F_A([a])=\exp(c_1(\mathcal P_a))=z$. Fourier inversion proves $o=[a]$. Finally $[a]=[b]$ in rational Chow groups if and only if their Albanese classes differ by torsion: the forward implication follows from Albanese, and the reverse from the identical Fourier exponentials of torsion-equivalent points. This proves the uniqueness assertion and the theorem.
\end{proof}

\begin{proposition}[A finite criterion for curves]
\label{c:curve-chow-criterion}
\status{Proved over every algebraically closed field}
Let $C$ be a smooth projective connected curve and $o\in\CH_0(C)_\Q$ have degree one. All point-and-diagonal cohomological relations hold in all its powers if and only if $M_o=0$ in $\CH^2(C^3)_\Q$. If $g\geq2$, this forces
\[
 o=\frac{K_C}{2g-2}\quad\text{in }\CH^1(C)_\Q.
\]
For $g=0$ or $1$ every degree-one rational zero-cycle has the required property. Over $\C$, the genus-at-least-two criterion is also the multiplicativity criterion in \citet[Propositions 3.1 and 3.5]{c_fulaterveervial2021}.
\end{proposition}

\begin{proof}
Necessity follows because $M_o$ is always cohomologically zero for a curve. Suppose it vanishes. Intersecting with $p_1^*o$ and pushing to the last two factors gives
\[
 \Delta_*o=o\boxtimes o.
\]
Intersecting instead with $\delta_{12}$ and pushing to the third factor gives
\[
 c_1(T_C)=(2-2g)o.
\]
In the second calculation, the small-diagonal term contributes $c_1(T_C)$, the term $\delta_{12}o_3$ contributes $(2-2g)o$, and the four nonzero remaining point terms cancel in pairs. These identities and the expansion of $M_o=0$ give all the local relations \eqref{c:pair-local} for $q_{ij}=\delta_{ij}-o_i-o_j$, with parameter $b=-2g$.

We check the remaining matching relation, including for a rational degree-one zero-cycle. Choose a base point $e$ and the Abel--Jacobi map $C\to J(C)$. The Jacobian's principal autoduality and the theorem of the square identify the pullback of its normalized symmetric Poincar\'e class, up to a common sign, with $\delta_{12}-e_1-e_2$. Translation of the Abel--Jacobi map by $z\in J(C)$ adds the corresponding algebraically trivial divisor to each factor. Since $J(C)(k)$ is divisible, we may choose $z$ so that the new pullback is $q_{12}$ in rational Chow groups, up to the same sign. Thus every matching monomial in the $q_{ij}$ is the pullback of a matching monomial in Poincar\'e divisors on a power of $J(C)$.

Let $\theta$ be a symmetric principal polarization on $J(C)$ and use integer parameters $n_1,\ldots,n_{2g+2}$. Pull back $\theta^{g+1}=0$ along the homomorphism $(x_i)\mapsto\sum n_ix_i$. The theorem of the square gives a polynomial identity whose quadratic form is
\[
 \sum_i n_i^2\theta_i+\sum_{i<j}n_in_jP_{ij}.
\]
The coefficient of $n_1\cdots n_{2g+2}$ in its $(g+1)$st power is $(g+1)!$ times the sum over all perfect matchings of the products of $P_{ij}$. A rational polynomial vanishing at every integer tuple has every coefficient zero. Pulling back to $C^{2g+2}$ proves precisely \eqref{c:sympl-pf}. Theorem~\ref{c:pair-ranks} now proves sufficiency in every power. These are algebraic arguments valid also in characteristic two.

For genus zero, $\CH_0(C)_\Q=\Q$. For genus one, choose an origin. The degree-zero Chow group is its Jacobian group, so divisibility shows that every degree-one rational zero-cycle equals a point class. Apply Theorem~\ref{c:abelian-point-classification}. The formula for $o$ when $g\geq2$ follows from the second displayed identity.
\end{proof}

\begin{proposition}[The exact additional obstruction for regular surfaces]
\label{c:surface-chow-criterion}
\status{Proved equivalence; does not assume finite-dimensionality}
Let $S/\C$ be a smooth projective surface with $b_1(S)=0$, and let $o$ have degree one. All the point-and-diagonal cohomological relations lift in all powers if and only if $M_o=0$ and $h(S)$ is Kimura finite-dimensional. If they hold, $o$ is uniquely determined: for every ample divisor $H$,
\[
 o=\frac{H^2}{\deg H^2}.
\]
\end{proposition}

\begin{proof}
The projector $\pi_2=\Delta-o\times S-S\times o$ has realization $H^2(S)$; the other two summands are Tate. Write $b=b_2(S)$. The alternating cross-matching determinant on $2(b+1)$ factors is $(b+1)!$ times the alternating projector on $h^2(S)^{\otimes(b+1)}$. Therefore the lifting of all relations implies $\bigwedge^{b+1}h^2(S)=0$, hence finite-dimensionality. It also implies $M_o=0$, since three positive-degree classes have zero product on such a surface.

Conversely, let $M_o=0$. Acting on $o$, on the diagonal, and on $H\boxtimes H$ gives respectively
\[
 \Delta_*o=o\boxtimes o,\qquad c_2(T_S)=(b+2)o,
 \qquad H^2=(\deg H^2)o.
\]
The same expansions as for the curve give all local relations \eqref{c:pair-local}, now with $b=b_2(S)$. Finite-dimensionality supplies the remaining alternating determinant: the realization is purely even of dimension $b$, so its Kimura exterior dimension is $b$. Theorem~\ref{c:pair-ranks} says that the local identities and these determinants generate every cohomological relation. This proves sufficiency and uniqueness. This formal argument is consistent with the larger divisor-generated criterion in \citet[Proposition 4.5]{c_fulaterveervial2021}; it explicitly applies to the narrower ring asked for here.
\end{proof}
